% Options for packages loaded elsewhere
% Options for packages loaded elsewhere
\PassOptionsToPackage{unicode}{hyperref}
\PassOptionsToPackage{hyphens}{url}
\PassOptionsToPackage{dvipsnames,svgnames,x11names}{xcolor}
%
\documentclass[
  letterpaper,
  DIV=11,
  numbers=noendperiod]{scrartcl}
\usepackage{xcolor}
\usepackage{amsmath,amssymb}
\setcounter{secnumdepth}{-\maxdimen} % remove section numbering
\usepackage{iftex}
\ifPDFTeX
  \usepackage[T1]{fontenc}
  \usepackage[utf8]{inputenc}
  \usepackage{textcomp} % provide euro and other symbols
\else % if luatex or xetex
  \usepackage{unicode-math} % this also loads fontspec
  \defaultfontfeatures{Scale=MatchLowercase}
  \defaultfontfeatures[\rmfamily]{Ligatures=TeX,Scale=1}
\fi
\usepackage{lmodern}
\ifPDFTeX\else
  % xetex/luatex font selection
\fi
% Use upquote if available, for straight quotes in verbatim environments
\IfFileExists{upquote.sty}{\usepackage{upquote}}{}
\IfFileExists{microtype.sty}{% use microtype if available
  \usepackage[]{microtype}
  \UseMicrotypeSet[protrusion]{basicmath} % disable protrusion for tt fonts
}{}
\makeatletter
\@ifundefined{KOMAClassName}{% if non-KOMA class
  \IfFileExists{parskip.sty}{%
    \usepackage{parskip}
  }{% else
    \setlength{\parindent}{0pt}
    \setlength{\parskip}{6pt plus 2pt minus 1pt}}
}{% if KOMA class
  \KOMAoptions{parskip=half}}
\makeatother
% Make \paragraph and \subparagraph free-standing
\makeatletter
\ifx\paragraph\undefined\else
  \let\oldparagraph\paragraph
  \renewcommand{\paragraph}{
    \@ifstar
      \xxxParagraphStar
      \xxxParagraphNoStar
  }
  \newcommand{\xxxParagraphStar}[1]{\oldparagraph*{#1}\mbox{}}
  \newcommand{\xxxParagraphNoStar}[1]{\oldparagraph{#1}\mbox{}}
\fi
\ifx\subparagraph\undefined\else
  \let\oldsubparagraph\subparagraph
  \renewcommand{\subparagraph}{
    \@ifstar
      \xxxSubParagraphStar
      \xxxSubParagraphNoStar
  }
  \newcommand{\xxxSubParagraphStar}[1]{\oldsubparagraph*{#1}\mbox{}}
  \newcommand{\xxxSubParagraphNoStar}[1]{\oldsubparagraph{#1}\mbox{}}
\fi
\makeatother

\usepackage{color}
\usepackage{fancyvrb}
\newcommand{\VerbBar}{|}
\newcommand{\VERB}{\Verb[commandchars=\\\{\}]}
\DefineVerbatimEnvironment{Highlighting}{Verbatim}{commandchars=\\\{\}}
% Add ',fontsize=\small' for more characters per line
\usepackage{framed}
\definecolor{shadecolor}{RGB}{241,243,245}
\newenvironment{Shaded}{\begin{snugshade}}{\end{snugshade}}
\newcommand{\AlertTok}[1]{\textcolor[rgb]{0.68,0.00,0.00}{#1}}
\newcommand{\AnnotationTok}[1]{\textcolor[rgb]{0.37,0.37,0.37}{#1}}
\newcommand{\AttributeTok}[1]{\textcolor[rgb]{0.40,0.45,0.13}{#1}}
\newcommand{\BaseNTok}[1]{\textcolor[rgb]{0.68,0.00,0.00}{#1}}
\newcommand{\BuiltInTok}[1]{\textcolor[rgb]{0.00,0.23,0.31}{#1}}
\newcommand{\CharTok}[1]{\textcolor[rgb]{0.13,0.47,0.30}{#1}}
\newcommand{\CommentTok}[1]{\textcolor[rgb]{0.37,0.37,0.37}{#1}}
\newcommand{\CommentVarTok}[1]{\textcolor[rgb]{0.37,0.37,0.37}{\textit{#1}}}
\newcommand{\ConstantTok}[1]{\textcolor[rgb]{0.56,0.35,0.01}{#1}}
\newcommand{\ControlFlowTok}[1]{\textcolor[rgb]{0.00,0.23,0.31}{\textbf{#1}}}
\newcommand{\DataTypeTok}[1]{\textcolor[rgb]{0.68,0.00,0.00}{#1}}
\newcommand{\DecValTok}[1]{\textcolor[rgb]{0.68,0.00,0.00}{#1}}
\newcommand{\DocumentationTok}[1]{\textcolor[rgb]{0.37,0.37,0.37}{\textit{#1}}}
\newcommand{\ErrorTok}[1]{\textcolor[rgb]{0.68,0.00,0.00}{#1}}
\newcommand{\ExtensionTok}[1]{\textcolor[rgb]{0.00,0.23,0.31}{#1}}
\newcommand{\FloatTok}[1]{\textcolor[rgb]{0.68,0.00,0.00}{#1}}
\newcommand{\FunctionTok}[1]{\textcolor[rgb]{0.28,0.35,0.67}{#1}}
\newcommand{\ImportTok}[1]{\textcolor[rgb]{0.00,0.46,0.62}{#1}}
\newcommand{\InformationTok}[1]{\textcolor[rgb]{0.37,0.37,0.37}{#1}}
\newcommand{\KeywordTok}[1]{\textcolor[rgb]{0.00,0.23,0.31}{\textbf{#1}}}
\newcommand{\NormalTok}[1]{\textcolor[rgb]{0.00,0.23,0.31}{#1}}
\newcommand{\OperatorTok}[1]{\textcolor[rgb]{0.37,0.37,0.37}{#1}}
\newcommand{\OtherTok}[1]{\textcolor[rgb]{0.00,0.23,0.31}{#1}}
\newcommand{\PreprocessorTok}[1]{\textcolor[rgb]{0.68,0.00,0.00}{#1}}
\newcommand{\RegionMarkerTok}[1]{\textcolor[rgb]{0.00,0.23,0.31}{#1}}
\newcommand{\SpecialCharTok}[1]{\textcolor[rgb]{0.37,0.37,0.37}{#1}}
\newcommand{\SpecialStringTok}[1]{\textcolor[rgb]{0.13,0.47,0.30}{#1}}
\newcommand{\StringTok}[1]{\textcolor[rgb]{0.13,0.47,0.30}{#1}}
\newcommand{\VariableTok}[1]{\textcolor[rgb]{0.07,0.07,0.07}{#1}}
\newcommand{\VerbatimStringTok}[1]{\textcolor[rgb]{0.13,0.47,0.30}{#1}}
\newcommand{\WarningTok}[1]{\textcolor[rgb]{0.37,0.37,0.37}{\textit{#1}}}

\usepackage{longtable,booktabs,array}
\newcounter{none} % for unnumbered tables
\usepackage{calc} % for calculating minipage widths
% Correct order of tables after \paragraph or \subparagraph
\usepackage{etoolbox}
\makeatletter
\patchcmd\longtable{\par}{\if@noskipsec\mbox{}\fi\par}{}{}
\makeatother
% Allow footnotes in longtable head/foot
\IfFileExists{footnotehyper.sty}{\usepackage{footnotehyper}}{\usepackage{footnote}}
\makesavenoteenv{longtable}
\usepackage{graphicx}
\makeatletter
\newsavebox\pandoc@box
\newcommand*\pandocbounded[1]{% scales image to fit in text height/width
  \sbox\pandoc@box{#1}%
  \Gscale@div\@tempa{\textheight}{\dimexpr\ht\pandoc@box+\dp\pandoc@box\relax}%
  \Gscale@div\@tempb{\linewidth}{\wd\pandoc@box}%
  \ifdim\@tempb\p@<\@tempa\p@\let\@tempa\@tempb\fi% select the smaller of both
  \ifdim\@tempa\p@<\p@\scalebox{\@tempa}{\usebox\pandoc@box}%
  \else\usebox{\pandoc@box}%
  \fi%
}
% Set default figure placement to htbp
\def\fps@figure{htbp}
\makeatother





\setlength{\emergencystretch}{3em} % prevent overfull lines

\providecommand{\tightlist}{%
  \setlength{\itemsep}{0pt}\setlength{\parskip}{0pt}}



 


\KOMAoption{captions}{tableheading}
\makeatletter
\@ifpackageloaded{caption}{}{\usepackage{caption}}
\AtBeginDocument{%
\ifdefined\contentsname
  \renewcommand*\contentsname{Table of contents}
\else
  \newcommand\contentsname{Table of contents}
\fi
\ifdefined\listfigurename
  \renewcommand*\listfigurename{List of Figures}
\else
  \newcommand\listfigurename{List of Figures}
\fi
\ifdefined\listtablename
  \renewcommand*\listtablename{List of Tables}
\else
  \newcommand\listtablename{List of Tables}
\fi
\ifdefined\figurename
  \renewcommand*\figurename{Figure}
\else
  \newcommand\figurename{Figure}
\fi
\ifdefined\tablename
  \renewcommand*\tablename{Table}
\else
  \newcommand\tablename{Table}
\fi
}
\@ifpackageloaded{float}{}{\usepackage{float}}
\floatstyle{ruled}
\@ifundefined{c@chapter}{\newfloat{codelisting}{h}{lop}}{\newfloat{codelisting}{h}{lop}[chapter]}
\floatname{codelisting}{Listing}
\newcommand*\listoflistings{\listof{codelisting}{List of Listings}}
\makeatother
\makeatletter
\makeatother
\makeatletter
\@ifpackageloaded{caption}{}{\usepackage{caption}}
\@ifpackageloaded{subcaption}{}{\usepackage{subcaption}}
\makeatother
\usepackage{bookmark}
\IfFileExists{xurl.sty}{\usepackage{xurl}}{} % add URL line breaks if available
\urlstyle{same}
\hypersetup{
  pdftitle={Week 05\_Gating sets},
  pdfauthor={Macarena López-Mayorga},
  colorlinks=true,
  linkcolor={blue},
  filecolor={Maroon},
  citecolor={Blue},
  urlcolor={Blue},
  pdfcreator={LaTeX via pandoc}}


\title{Week 05\_Gating sets}
\author{Macarena López-Mayorga}
\date{2026-04-23}
\begin{document}
\maketitle


\begin{Shaded}
\begin{Highlighting}[]
\NormalTok{Folder }\OtherTok{\textless{}{-}} \FunctionTok{file.path}\NormalTok{(}\StringTok{"C:/Users/Macarena/CytometryInR/course/05\_GatingSets/data"}\NormalTok{)}
\NormalTok{fcs\_files }\OtherTok{\textless{}{-}} \FunctionTok{list.files}\NormalTok{(Folder, }\AttributeTok{pattern=}\StringTok{".fcs"}\NormalTok{, }\AttributeTok{full.names=}\ConstantTok{TRUE}\NormalTok{)}
\NormalTok{fcs\_files}
\end{Highlighting}
\end{Shaded}

{[}1{]}
``C:/Users/Macarena/CytometryInR/course/05\_GatingSets/data/2025\_07\_26\_AB\_02\_INF052\_00\_Ctrl.fcs''
{[}2{]}
``C:/Users/Macarena/CytometryInR/course/05\_GatingSets/data/2025\_07\_26\_AB\_02\_INF052\_00\_SEB.fcs''
{[}3{]}
``C:/Users/Macarena/CytometryInR/course/05\_GatingSets/data/2025\_07\_26\_AB\_02\_INF100\_00\_Ctrl.fcs''
{[}4{]}
``C:/Users/Macarena/CytometryInR/course/05\_GatingSets/data/2025\_07\_26\_AB\_02\_INF100\_00\_SEB.fcs''
{[}5{]}
``C:/Users/Macarena/CytometryInR/course/05\_GatingSets/data/2025\_07\_26\_AB\_02\_INF179\_00\_Ctrl.fcs''
{[}6{]}
``C:/Users/Macarena/CytometryInR/course/05\_GatingSets/data/2025\_07\_26\_AB\_02\_INF179\_00\_SEB.fcs''

We are going to name one of them as ``fcs-files{[}x{]}''. In our case,
the first file

\begin{Shaded}
\begin{Highlighting}[]
\NormalTok{fcs\_files[}\DecValTok{1}\NormalTok{]}
\end{Highlighting}
\end{Shaded}

{[}1{]}
``C:/Users/Macarena/CytometryInR/course/05\_GatingSets/data/2025\_07\_26\_AB\_02\_INF052\_00\_Ctrl.fcs''

\begin{Shaded}
\begin{Highlighting}[]
\FunctionTok{library}\NormalTok{(flowCore)}
\end{Highlighting}
\end{Shaded}

\begin{Shaded}
\begin{Highlighting}[]
\NormalTok{flowFrame }\OtherTok{\textless{}{-}} \FunctionTok{read.FCS}\NormalTok{(}\AttributeTok{filename=}\NormalTok{fcs\_files[}\DecValTok{1}\NormalTok{], }\AttributeTok{truncate\_max\_range =} \ConstantTok{FALSE}\NormalTok{, }\AttributeTok{transformation =} \ConstantTok{FALSE}\NormalTok{)}
\NormalTok{flowFrame}
\end{Highlighting}
\end{Shaded}

flowFrame object `2025\_07\_26\_AB\_02\_INF052\_00\_Ctrl.fcs' with 10000
cells and 43 observables: name desc range minRange maxRange
\(P1            Time        NA    896745         0     878.809\)P2 SSC-W
NA 4194304 0 4194303.000
\(P3           SSC-H        NA   4194304         0 4194303.000\)P4 SSC-A
NA 4194304 0 4194303.000
\(P5           FSC-W        NA   4194304         0 4194303.000
...             ...       ...       ...       ...         ...\)P39
APC-R700-A CD107a 4194304 -111 4192506
\(P40   Zombie NIR-A Viability   4194304      -111     4192506\)P41
APC-Fire 750-A CD27 4194304 -111 4192506
\(P42 APC-Fire 810-A      CCR7   4194304      -111     4192506\)P43 AF-A
NA 4194304 -111 4194303 472 keywords are stored in the `description'
slot

\begin{Shaded}
\begin{Highlighting}[]
\FunctionTok{str}\NormalTok{(fcs\_files)}
\end{Highlighting}
\end{Shaded}

chr {[}1:6{]}
``C:/Users/Macarena/CytometryInR/course/05\_GatingSets/data/2025\_07\_26\_AB\_02\_INF052\_00\_Ctrl.fcs''
\ldots{}

All of them are character class.

If we want to read multiple fcs files at once, we will need another
function:

::: \{.cell\}

\begin{Shaded}
\begin{Highlighting}[]
\NormalTok{flowSet }\OtherTok{\textless{}{-}} \FunctionTok{read.flowSet}\NormalTok{(}\AttributeTok{files=}\NormalTok{fcs\_files, }\AttributeTok{truncate\_max\_range =} \ConstantTok{FALSE}\NormalTok{, }\AttributeTok{transformation =} \ConstantTok{FALSE}\NormalTok{)}
\NormalTok{flowSet}
\end{Highlighting}
\end{Shaded}

:::

A flowSet with 6 experiments.

column names(43): Time SSC-W \ldots{} APC-Fire 810-A AF-A

Also, we can select those fcs files that we can specifically read:

\begin{Shaded}
\begin{Highlighting}[]
\FunctionTok{read.flowSet}\NormalTok{(}\AttributeTok{files=}\NormalTok{fcs\_files[}\FunctionTok{c}\NormalTok{(}\DecValTok{1}\NormalTok{, }\DecValTok{3}\SpecialCharTok{:}\DecValTok{4}\NormalTok{)], }\AttributeTok{truncate\_max\_range =} \ConstantTok{FALSE}\NormalTok{, }\AttributeTok{transformation =} \ConstantTok{FALSE}\NormalTok{)}
\end{Highlighting}
\end{Shaded}

A flowSet with 3 experiments.

column names(43): Time SSC-W \ldots{} APC-Fire 810-A AF-A

we can see that read.flowSet() has created a ``flowSet'' class object.

\begin{Shaded}
\begin{Highlighting}[]
\FunctionTok{class}\NormalTok{(flowSet)}
\end{Highlighting}
\end{Shaded}

{[}1{]} ``flowSet'' attr(,``package'') {[}1{]} ``flowCore''

In the right secondary side bar, we can see that ``flowSet'' is another
Bioconductor style S4-type object that within its frame slot contains
individual ``flowFrames''. If instead of class() we had used str(), we
would have seen a similar output ton what we see in the Variables panel.

\begin{Shaded}
\begin{Highlighting}[]
\FunctionTok{str}\NormalTok{(flowSet)}
\end{Highlighting}
\end{Shaded}

Formal class `flowSet' {[}package ``flowCore''{]} with 2 slots ..@
frames :\textless environment: 0x000001fb5d4cf5a8\textgreater{} ..@
phenoData:Formal class `AnnotatedDataFrame' {[}package ``Biobase''{]}
with 4 slots .. .. ..@ varMetadata :`data.frame': 1 obs. of 1 variable:
.. .. .. ..\$ labelDescription: chr ``Name'' .. .. ..@ data
:`data.frame': 6 obs. of 1 variable: .. .. .. ..\$ name: `AsIs' chr
{[}1:6{]} ``2025\_07\_26\_AB\_02\_INF052\_00\_Ctrl.fcs''
``2025\_07\_26\_AB\_02\_INF052\_00\_SEB.fcs''
``2025\_07\_26\_AB\_02\_INF100\_00\_Ctrl.fcs''
``2025\_07\_26\_AB\_02\_INF100\_00\_SEB.fcs'' \ldots{} .. .. ..@
dimLabels : chr {[}1:2{]} ``rowNames'' ``columnNames'' .. .. ..@
.\_\_classVersion\_\_:Formal class `Versions' {[}package ``Biobase''{]}
with 1 slot .. .. .. .. ..@ .Data:List of 1 .. .. .. .. .. ..\$ : int
{[}1:3{]} 1 1 0 .. .. .. .. ..\$ names: chr ``AnnotatedDataFrame''

REMINDER: While not today's focus, remember we could access individual
components inside the flowSet using the @ accessors covered during
Week3.

One consideration of this method is the contents of your .fcs files are
read into your computer's random access memory (RAM). While for
individual .fcs files or small experiments this will not present a
problem for most modern computers, when working with large spectral flow
cytometry files containing millions of events (or trying to analyze many
.fcs files at once), you may encounter situations where you can quickly
exceed your computers available RAM. To build some contextual
understanding of the problem, let's learn how to check how much memory
is being used by our individual variables/objects within our R session.
We will primarily use the lobstr R packages obj\_size() function, as it
better handles evaluating complicated objects than base R's
object.size() function.

::: \{.cell\}

\begin{Shaded}
\begin{Highlighting}[]
\FunctionTok{object.size}\NormalTok{(flowFrame)}
\end{Highlighting}
\end{Shaded}

:::

3570976 bytes

::: \{.cell\}

\begin{Shaded}
\begin{Highlighting}[]
\FunctionTok{install.packages}\NormalTok{(}\StringTok{"lobstr"}\NormalTok{)}
\end{Highlighting}
\end{Shaded}

:::

::: \{.cell\}

\begin{Shaded}
\begin{Highlighting}[]
 \FunctionTok{library}\NormalTok{(lobstr)}
\FunctionTok{obj\_size}\NormalTok{(flowFrame) }
\end{Highlighting}
\end{Shaded}

:::

3.53 MB FlowSet object contains multiples flowFrame objects, so the size
should be bigger:

::: \{.cell\}

\begin{Shaded}
\begin{Highlighting}[]
\FunctionTok{obj\_size}\NormalTok{(flowSet)}
\end{Highlighting}
\end{Shaded}

:::

20.99 MB

If we were curious how much memory total we are using within R at the
current moment, we can check using the mem\_used() function:

::: \{.cell\}

\begin{Shaded}
\begin{Highlighting}[]
\FunctionTok{mem\_used}\NormalTok{()}
\end{Highlighting}
\end{Shaded}

:::

141.95 MB

Ultimately, how many .fcs files you are able to read in and interact
with before running out of available RAM memory space will be dictated
by your individual computers hardware configuration. There are various
ways you can check programmatically how much RAM your computer has
available, although the specific functions will vary depending on your
computers operating system, since they often involve system-level code
outside R. Using the ps R package's ps\_system\_memory() function is one
of the easier ways for Windows users.

To simplify the process, here is an additional example of where a
conditional can prove useful, allowing us to check in an operating
system specific manner. It takes the output of the Sys.info() function,
namely the ``sysname'' argument and then retrieves the relavent
function.

\begin{Shaded}
\begin{Highlighting}[]
\NormalTok{OperatingSystem }\OtherTok{\textless{}{-}} \FunctionTok{Sys.info}\NormalTok{()[[}\StringTok{"sysname"}\NormalTok{]]}
\ControlFlowTok{if}\NormalTok{ (OperatingSystem }\SpecialCharTok{==} \StringTok{"Windows"}\NormalTok{)}
\NormalTok{Memory }\OtherTok{\textless{}{-}}\NormalTok{ ps}\SpecialCharTok{::}\FunctionTok{ps\_system\_memory}\NormalTok{()}
\FunctionTok{message}\NormalTok{(}\StringTok{"Total GB "}\NormalTok{, }\FunctionTok{round}\NormalTok{(Memory}\SpecialCharTok{$}\NormalTok{total }\SpecialCharTok{/} \DecValTok{1024}\SpecialCharTok{\^{}}\DecValTok{3}\NormalTok{, }\DecValTok{2}\NormalTok{))}
\FunctionTok{message}\NormalTok{(}\StringTok{"Free GB "}\NormalTok{, }\FunctionTok{round}\NormalTok{(Memory}\SpecialCharTok{$}\NormalTok{free }\SpecialCharTok{/} \DecValTok{1024}\SpecialCharTok{\^{}}\DecValTok{3}\NormalTok{, }\DecValTok{2}\NormalTok{))}
\end{Highlighting}
\end{Shaded}

Total GB 15.78 Free GB 7.69

\begin{Shaded}
\begin{Highlighting}[]
\FunctionTok{install.packages}\NormalTok{(}\StringTok{"ps"}\NormalTok{)}
\end{Highlighting}
\end{Shaded}

\begin{Shaded}
\begin{Highlighting}[]
\FunctionTok{library}\NormalTok{(ps)}
\NormalTok{Memory }\OtherTok{\textless{}{-}}\NormalTok{ ps}\SpecialCharTok{::}\FunctionTok{ps\_system\_memory}\NormalTok{()}
\FunctionTok{message}\NormalTok{(}\StringTok{"Total GB "}\NormalTok{, }\FunctionTok{round}\NormalTok{(Memory}\SpecialCharTok{$}\NormalTok{total }\SpecialCharTok{/} \DecValTok{1024}\SpecialCharTok{\^{}}\DecValTok{3}\NormalTok{, }\DecValTok{2}\NormalTok{))}
\end{Highlighting}
\end{Shaded}

Total GB 15.78

\begin{Shaded}
\begin{Highlighting}[]
\FunctionTok{message}\NormalTok{(}\StringTok{"Free GB "}\NormalTok{, }\FunctionTok{round}\NormalTok{(Memory}\SpecialCharTok{$}\NormalTok{free }\SpecialCharTok{/} \DecValTok{1024}\SpecialCharTok{\^{}}\DecValTok{3}\NormalTok{, }\DecValTok{2}\NormalTok{))}
\end{Highlighting}
\end{Shaded}

Free GB 7.83

In addition to the flowCore R package, additional flow cytometry
infrastructure support is provided by the flowWorkspace package. Instead
of the reading all the .fcs files contents into active RAM,
flowWorkspace reduces the memory overhead by using ``pointers'' to
interact with the object in it's current storage location (either on
your harddrive, SSD, etc.), only reading in components to RAM as needed.

\begin{Shaded}
\begin{Highlighting}[]
\FunctionTok{library}\NormalTok{(flowWorkspace)}
\end{Highlighting}
\end{Shaded}

As part of improvements to flowWorkspace, some behavior of GatingSet
objects has changed. For details, please read the section titled ``The
cytoframe and cytoset classes'' in the package vignette:

vignette(``flowWorkspace-Introduction'', ``flowWorkspace'')

::: \{.cell\}

\begin{Shaded}
\begin{Highlighting}[]
\NormalTok{cytoframe }\OtherTok{\textless{}{-}} \FunctionTok{load\_cytoframe\_from\_fcs}\NormalTok{(fcs\_files[}\DecValTok{1}\NormalTok{], }\AttributeTok{truncate\_max\_range =} \ConstantTok{FALSE}\NormalTok{, }\AttributeTok{transformation =} \ConstantTok{FALSE}\NormalTok{)}
\NormalTok{cytoframe}
\end{Highlighting}
\end{Shaded}

:::

cytoframe object `2025\_07\_26\_AB\_02\_INF052\_00\_Ctrl.fcs' with 10000
cells and 43 observables: name desc range minRange maxRange
\(P1            Time        NA     878.809         0     878.809\)P2
SSC-W NA 4194303.000 0 4194303.000
\(P3           SSC-H        NA 4194303.000         0 4194303.000\)P4
SSC-A NA 4194303.000 0 4194303.000
\(P5           FSC-W        NA 4194303.000         0 4194303.000
...             ...       ...         ...       ...         ...\)P39
APC-R700-A CD107a 4192506 -111 4192506
\(P40   Zombie NIR-A Viability     4192506      -111     4192506\)P41
APC-Fire 750-A CD27 4192506 -111 4192506
\(P42 APC-Fire 810-A      CCR7     4192506      -111     4192506\)P43
AF-A NA 4194303 -111 4194303 472 keywords are stored in the
`description' slot row names(0):

\begin{Shaded}
\begin{Highlighting}[]
\FunctionTok{class}\NormalTok{(cytoframe)}
\end{Highlighting}
\end{Shaded}

{[}1{]} ``cytoframe'' attr(,``package'') {[}1{]} ``flowWorkspace''

Which also still errors out when not given a scalar object

\begin{Shaded}
\begin{Highlighting}[]
\FunctionTok{load\_cytoframe\_from\_fcs}\NormalTok{(fcs\_files, }\AttributeTok{truncate\_max\_range =} \ConstantTok{FALSE}\NormalTok{, }\AttributeTok{transformation =} \ConstantTok{FALSE}\NormalTok{)}
\end{Highlighting}
\end{Shaded}

Error: ! Expected string vector of length 1

We can use Cytoset to handle multiples fcs files:

::: \{.cell\}

\begin{Shaded}
\begin{Highlighting}[]
\NormalTok{cytoset }\OtherTok{\textless{}{-}} \FunctionTok{load\_cytoset\_from\_fcs}\NormalTok{(fcs\_files, }\AttributeTok{truncate\_max\_range =} \ConstantTok{FALSE}\NormalTok{, }\AttributeTok{transformation =} \ConstantTok{FALSE}\NormalTok{)}
\NormalTok{cytoset}
\end{Highlighting}
\end{Shaded}

:::

A cytoset with 6 samples.

column names: Time, SSC-W, SSC-H, SSC-A, FSC-W, FSC-H, FSC-A, SSC-B-W,
SSC-B-H, SSC-B-A, BUV395-A, BUV563-A, BUV615-A, BUV661-A, BUV737-A,
BUV805-A, Pacific Blue-A, BV480-A, BV570-A, BV605-A, BV650-A, BV711-A,
BV750-A, BV786-A, Alexa Fluor 488-A, Spark Blue 550-A, Spark Blue 574-A,
RB613-A, RB705-A, RB780-A, PE-A, PE-Dazzle594-A, PE-Cy5-A, PE-Fire
700-A, PE-Fire 744-A, PE-Vio770-A, APC-A, Alexa Fluor 647-A, APC-R700-A,
Zombie NIR-A, APC-Fire 750-A, APC-Fire 810-A, AF-A

\begin{Shaded}
\begin{Highlighting}[]
\FunctionTok{class}\NormalTok{(cytoset)}
\end{Highlighting}
\end{Shaded}

{[}1{]} ``cytoset'' attr(,``package'') {[}1{]} ``flowWorkspace''

Unlike ``flowFrame'' and ``flowSet'', when we run str(), for
``cytoframe'' and ``cytoset'' objects we don't get back quite as much
information.

\begin{Shaded}
\begin{Highlighting}[]
\FunctionTok{str}\NormalTok{(cytoframe)}
\end{Highlighting}
\end{Shaded}

Formal class `cytoframe' {[}package ``flowWorkspace''{]} with 5 slots
..@ pointer : ..@ use.exprs : logi TRUE ..@ exprs : num{[}0 , 0 {]} ..@
parameters :Formal class `AnnotatedDataFrame' {[}package ``Biobase''{]}
with 4 slots .. .. ..@ varMetadata :`data.frame': 5 obs. of 1 variable:
.. .. .. ..\$ labelDescription: chr {[}1:5{]} ``Name of Parameter''
``Description of Parameter'' ``Range of Parameter'' ``Minimum Parameter
Value after Transformation'' \ldots{} .. .. ..@ data :`data.frame': 0
obs. of 5 variables: .. .. .. ..\$ name : chr(0) .. .. .. ..\$
description: chr(0) .. .. .. ..\$ range : num(0) .. .. .. ..\$ minRange
: num(0) .. .. .. ..\$ maxRange : num(0) .. .. ..@ dimLabels : chr
{[}1:2{]} ``rowNames'' ``columnNames'' .. .. ..@
.\_\_classVersion\_\_:Formal class `Versions' {[}package ``Biobase''{]}
with 1 slot .. .. .. .. ..@ .Data:List of 1 .. .. .. .. .. ..\$ : int
{[}1:3{]} 1 1 0 .. .. .. .. ..\$ names: chr ``AnnotatedDataFrame'' ..@
description:List of 1 .. ..\$ note: chr ``empty''

::: \{.cell\}

\begin{Shaded}
\begin{Highlighting}[]
\FunctionTok{str}\NormalTok{(cytoset)}
\end{Highlighting}
\end{Shaded}

:::

Formal class `cytoset' {[}package ``flowWorkspace''{]} with 3 slots ..@
pointer : ..@ frames :\textless environment:
0x000001e3cf33dee8\textgreater{} ..@ phenoData:Formal class
`AnnotatedDataFrame' {[}package ``Biobase''{]} with 4 slots .. .. ..@
varMetadata :`data.frame': 0 obs. of 1 variable: .. .. .. ..\$
labelDescription: chr(0) .. .. ..@ data :`data.frame': 0 obs. of 0
variables .. .. ..@ dimLabels : chr {[}1:2{]} ``rowNames''
``columnNames'' .. .. ..@ .\_\_classVersion\_\_:Formal class `Versions'
{[}package ``Biobase''{]} with 1 slot .. .. .. .. ..@ .Data:List of 1 ..
.. .. .. .. ..\$ : int {[}1:3{]} 1 1 0 .. .. .. .. ..\$ names: chr
``AnnotatedDataFrame''

This is similarly the case when glancing at the right secondary side
bar, as the respective objects under variables appear to have empty
matrices where normally we would have seen the MFI values. Due to
flowWorkspace use of pointers, the missing data remains stored on the
drive, only being retrieved right before it is required. This reduces
the overall RAM utilization. Let's double check the differences in
memory utilization for flowFrame/cytoframe:

::: \{.cell\}

\begin{Shaded}
\begin{Highlighting}[]
\FunctionTok{obj\_size}\NormalTok{(flowFrame)}
\end{Highlighting}
\end{Shaded}

:::

3.53 MB

\begin{Shaded}
\begin{Highlighting}[]
\FunctionTok{obj\_size}\NormalTok{(cytoframe)}
\end{Highlighting}
\end{Shaded}

5.40 kB

And similarly the case for flowSet and cytoset:

\begin{Shaded}
\begin{Highlighting}[]
\FunctionTok{obj\_size}\NormalTok{(flowSet)}
\end{Highlighting}
\end{Shaded}

20.99 MB

\begin{Shaded}
\begin{Highlighting}[]
\FunctionTok{obj\_size}\NormalTok{(cytoset)}
\end{Highlighting}
\end{Shaded}

3.88 kB

Despite both R packages having been around for a while, many
Bioconductor and GitHub often only implement methods to handle either
flowFrames or cytoframes (although newer R packages are now allowing for
both). Consequently, as we move forward in the course, it helps to be
aware of which ones we are working with, and have the ability to
interconvert between them as needed.

To go from a flowFrame to a cytoframe, we can use the
flowFrame\_to\_cytoframe() function

\begin{Shaded}
\begin{Highlighting}[]
\NormalTok{ConvertedToCytoframe }\OtherTok{\textless{}{-}} \FunctionTok{flowFrame\_to\_cytoframe}\NormalTok{(flowFrame)}
\NormalTok{ConvertedToCytoframe}
\end{Highlighting}
\end{Shaded}

cytoframe object `file207836ed2889' with 10000 cells and 43 observables:
name desc range minRange maxRange
\(P1            Time        NA     878.809         0     878.809\)P2
SSC-W NA 4194303.000 0 4194303.000
\(P3           SSC-H        NA 4194303.000         0 4194303.000\)P4
SSC-A NA 4194303.000 0 4194303.000
\(P5           FSC-W        NA 4194303.000         0 4194303.000
...             ...       ...         ...       ...         ...\)P39
APC-R700-A CD107a 4192506 -111 4192506
\(P40   Zombie NIR-A Viability     4192506      -111     4192506\)P41
APC-Fire 750-A CD27 4192506 -111 4192506
\(P42 APC-Fire 810-A      CCR7     4192506      -111     4192506\)P43
AF-A NA 4194303 -111 4194303 472 keywords are stored in the
`description' slot row names(0):

\begin{Shaded}
\begin{Highlighting}[]
\FunctionTok{obj\_size}\NormalTok{(ConvertedToCytoframe)}
\end{Highlighting}
\end{Shaded}

5.40 kB

To go from a cytoframe to a flowFrame, we can use the
cytoframe\_to\_flowFrame() function

\begin{Shaded}
\begin{Highlighting}[]
\NormalTok{ConvertedToFlowframe }\OtherTok{\textless{}{-}}\NormalTok{ flowWorkspace}\SpecialCharTok{::}\FunctionTok{cytoframe\_to\_flowFrame}\NormalTok{(cytoframe)}
\NormalTok{ConvertedToFlowframe}
\end{Highlighting}
\end{Shaded}

flowFrame object `2025\_07\_26\_AB\_02\_INF052\_00\_Ctrl.fcs' with 10000
cells and 43 observables: name desc range minRange maxRange
\(P1            Time        NA     878.809         0     878.809\)P2
SSC-W NA 4194303.000 0 4194303.000
\(P3           SSC-H        NA 4194303.000         0 4194303.000\)P4
SSC-A NA 4194303.000 0 4194303.000
\(P5           FSC-W        NA 4194303.000         0 4194303.000
...             ...       ...         ...       ...         ...\)P39
APC-R700-A CD107a 4192506 -111 4192506
\(P40   Zombie NIR-A Viability     4192506      -111     4192506\)P41
APC-Fire 750-A CD27 4192506 -111 4192506
\(P42 APC-Fire 810-A      CCR7     4192506      -111     4192506\)P43
AF-A NA 4194303 -111 4194303 472 keywords are stored in the
`description' slot

This other code gives the same result:

\begin{Shaded}
\begin{Highlighting}[]
\NormalTok{ConvertedToFlowframe }\OtherTok{\textless{}{-}} \FunctionTok{cytoframe\_to\_flowFrame}\NormalTok{(cytoframe)}
\NormalTok{ConvertedToFlowframe}
\end{Highlighting}
\end{Shaded}

\begin{Shaded}
\begin{Highlighting}[]
\FunctionTok{obj\_size}\NormalTok{(ConvertedToFlowframe)}
\end{Highlighting}
\end{Shaded}

3.53 MB

To go from a flowSet to a cytoSet, we can use the flowSet\_to\_cytoset()
function:

\begin{Shaded}
\begin{Highlighting}[]
\NormalTok{ConvertedToCytoset }\OtherTok{\textless{}{-}} \FunctionTok{flowSet\_to\_cytoset}\NormalTok{(flowSet)}
\NormalTok{ConvertedToCytoset}
\end{Highlighting}
\end{Shaded}

A cytoset with 6 samples.

column names: Time, SSC-W, SSC-H, SSC-A, FSC-W, FSC-H, FSC-A, SSC-B-W,
SSC-B-H, SSC-B-A, BUV395-A, BUV563-A, BUV615-A, BUV661-A, BUV737-A,
BUV805-A, Pacific Blue-A, BV480-A, BV570-A, BV605-A, BV650-A, BV711-A,
BV750-A, BV786-A, Alexa Fluor 488-A, Spark Blue 550-A, Spark Blue 574-A,
RB613-A, RB705-A, RB780-A, PE-A, PE-Dazzle594-A, PE-Cy5-A, PE-Fire
700-A, PE-Fire 744-A, PE-Vio770-A, APC-A, Alexa Fluor 647-A, APC-R700-A,
Zombie NIR-A, APC-Fire 750-A, APC-Fire 810-A, AF-A

\begin{Shaded}
\begin{Highlighting}[]
\FunctionTok{obj\_size}\NormalTok{(ConvertedToCytoset)}
\end{Highlighting}
\end{Shaded}

3.88 kB

\begin{Shaded}
\begin{Highlighting}[]
\NormalTok{ConvertedToFlowset }\OtherTok{\textless{}{-}} \FunctionTok{cytoset\_to\_flowSet}\NormalTok{(flowSet)}
\NormalTok{ConvertedToFlowset}
\end{Highlighting}
\end{Shaded}

A flowSet with 6 experiments.

column names(43): Time SSC-W \ldots{} APC-Fire 810-A AF-A

\begin{Shaded}
\begin{Highlighting}[]
\FunctionTok{obj\_size}\NormalTok{(ConvertedToFlowset)}
\end{Highlighting}
\end{Shaded}

20.99 MB

Fortunately, regardless of whether we are using flowFrame/flowSet (RAM)
and cytoframe/cytoset (memory pointers), both routes end up converging
at the next step, where the underlying .fcs files are passed off to the
GatingSet() function.

\begin{Shaded}
\begin{Highlighting}[]
\NormalTok{GatingSet1 }\OtherTok{\textless{}{-}} \FunctionTok{GatingSet}\NormalTok{(flowSet)}
\NormalTok{GatingSet1 }
\end{Highlighting}
\end{Shaded}

A GatingSet with 6 samples

\begin{Shaded}
\begin{Highlighting}[]
\FunctionTok{class}\NormalTok{(GatingSet1)}
\end{Highlighting}
\end{Shaded}

{[}1{]} ``GatingSet'' attr(,``package'') {[}1{]} ``flowWorkspace''

\begin{Shaded}
\begin{Highlighting}[]
\NormalTok{GatingSet2 }\OtherTok{\textless{}{-}} \FunctionTok{GatingSet}\NormalTok{(cytoset)}
\NormalTok{GatingSet2}
\end{Highlighting}
\end{Shaded}

A GatingSet with 6 samples

\begin{Shaded}
\begin{Highlighting}[]
\FunctionTok{class}\NormalTok{(GatingSet2)}
\end{Highlighting}
\end{Shaded}

{[}1{]} ``GatingSet'' attr(,``package'') {[}1{]} ``flowWorkspace''

As we prefaced in the background, beyond the .fcs files themselves, we
need infrastructural elements with which to interact with the underlying
data, which allows us to organize the various files, transform (scale),
compensate (for conventional flow), visualize, derrive statistics, etc.
A GatingSet serves as the infrastructural framework that allows us to do
this in R.

If we investigate our current GatingSet objects, we won't see much. This
will change as we start layering on additional elements. However, rather
than try to cram everything into a single week, we will explore in
greater depth the individual components over the next three weeks.
Instead, for the rest of today, we will work backward, by exploring a
GatingSet objecct and what it is capable of doing once fully assembled.

The CytoML R package (also maintained by Mike Jiang) is a sister package
to the flowWorkspace. It's main purpose is to permit bringing in
existing FlowJo, Diva and Cytobank Workspaces, with all their gates,
transformations, etc. into R as fully assembled GatingSet objects. For
those who already use one of these commercial softwares, it can be quite
useful tool.

Since our goal is to examine a fully assembled GatingSet object, we will
be using it today to bring in a FlowJo workspace to R. However, since
this is a free Cytometry in R course, and not about to have everyone pay
for a license for a one-off topic, in the pre-course Floreada
walkthrough I documented how to convert a free Floreada.io workspae into
a FlowJo.wsp that can also be used (please note that as of early 2026,
some scaling bugs may be present and require troubleshooting).

To get started, let's first attach CytoML to our local environment via
the library() call.

\begin{Shaded}
\begin{Highlighting}[]
\FunctionTok{library}\NormalTok{(CytoML)}
\end{Highlighting}
\end{Shaded}

The .wsp files within this week's data where created via Floreada.io.
The main difference between the two files is one is a copy of the
original that was opened within FlowJo, and subsequently swtiched from
logicle to bi-exponential transformation.

We will need to provide the appropiate file path for our desired .wsp
file. We can start by identifying which are present using list.files()

\begin{Shaded}
\begin{Highlighting}[]
\NormalTok{Folder }\OtherTok{\textless{}{-}} \FunctionTok{file.path}\NormalTok{(}\StringTok{"C:/Users/Macarena/CytometryInR/course/05\_GatingSets/data"}\NormalTok{)}
\NormalTok{FlowJoWsp }\OtherTok{\textless{}{-}} \FunctionTok{list.files}\NormalTok{(}\AttributeTok{path =}\NormalTok{ Folder, }\AttributeTok{pattern =} \StringTok{".wsp"}\NormalTok{, }\AttributeTok{full =} \ConstantTok{TRUE}\NormalTok{)}
\NormalTok{FlowJoWsp}
\end{Highlighting}
\end{Shaded}

{[}1{]}
``C:/Users/Macarena/CytometryInR/course/05\_GatingSets/data/FlowJoWSP\_OpenedCopy.wsp''
{[}2{]}
``C:/Users/Macarena/CytometryInR/course/05\_GatingSets/data/FlowJoWSP\_Unopened.wsp''

In our case, we will proceed by using str\_detect() to select the .wsp
that contains the pattern ``Opened''

\begin{Shaded}
\begin{Highlighting}[]
\NormalTok{ThisWorkspace }\OtherTok{\textless{}{-}}\NormalTok{ FlowJoWsp[stringr}\SpecialCharTok{::}\FunctionTok{str\_detect}\NormalTok{(FlowJoWsp, }\StringTok{"Opened"}\NormalTok{)]}
\NormalTok{ThisWorkspace}
\end{Highlighting}
\end{Shaded}

{[}1{]}
``C:/Users/Macarena/CytometryInR/course/05\_GatingSets/data/FlowJoWSP\_OpenedCopy.wsp''

With our single .wsp filepath now identified, we can now proceed to set
up the intermediate object using open\_flowjo\_xml()

\begin{Shaded}
\begin{Highlighting}[]
\NormalTok{ws }\OtherTok{\textless{}{-}} \FunctionTok{open\_flowjo\_xml}\NormalTok{(ThisWorkspace)}
\NormalTok{ws}
\end{Highlighting}
\end{Shaded}

File location:
C:/Users/Macarena/CytometryInR/course/05\_GatingSets/data/FlowJoWSP\_OpenedCopy.wsp

Groups in Workspace Name Num.Samples 1 All Samples 6

\begin{Shaded}
\begin{Highlighting}[]
\FunctionTok{class}\NormalTok{(ws)}
\end{Highlighting}
\end{Shaded}

{[}1{]} ``flowjo\_workspace'' attr(,``package'') {[}1{]} ``CytoML''

Having set up the intermediate flowjo\_workspace object, we can attempt
to read in the actual data from the .wsp into a GatingSet using the
flowjo\_to\_gatingset() function.

However, due to how I named the original .fcs files (``GROUPNAME'' being
individual specimens, ``TUBENAME'' being either Ctrl or SEB), and
downsampled to the same number of cells, we will encounter the following
error:

\begin{Shaded}
\begin{Highlighting}[]
\NormalTok{gs }\OtherTok{\textless{}{-}} \FunctionTok{flowjo\_to\_gatingset}\NormalTok{(}\AttributeTok{ws=}\NormalTok{ws, }\AttributeTok{name=}\DecValTok{1}\NormalTok{, }\AttributeTok{path =}\NormalTok{ Folder)}
\end{Highlighting}
\end{Shaded}

Error: ! Multiple FCS files match sample 00\_Ctrl.fcs by filename, event
count, and keywords. Candidates are:
C:\Users\Macarena\CytometryInR\course\textbackslash05\_GatingSets\data\textbackslash2025\_07\_26\_AB\_02\_INF052\_00\_Ctrl.fcs
C:\Users\Macarena\CytometryInR\course\textbackslash05\_GatingSets\data\textbackslash2025\_07\_26\_AB\_02\_INF100\_00\_Ctrl.fcs
C:\Users\Macarena\CytometryInR\course\textbackslash05\_GatingSets\data\textbackslash2025\_07\_26\_AB\_02\_INF179\_00\_Ctrl.fcs
Please move incorrect files out of this directory or its subdirectories.
Show Traceback

\begin{Shaded}
\begin{Highlighting}[]
\NormalTok{gs}
\end{Highlighting}
\end{Shaded}

Error: ! object `gs' not found

As with any error, my first move is to check the help documentation. In
this case, my initial response is to see if I can identify an argument
that will help differentiate between the names for each specimen.

\begin{Shaded}
\begin{Highlighting}[]
\NormalTok{?flowjo\_to\_gatingset}
\end{Highlighting}
\end{Shaded}

parseWorkspace(obj, \ldots)

\subsection{S4 method for signature
`flowjo\_workspace'}\label{s4-method-for-signature-flowjo_workspace}

parseWorkspace(obj, \ldots)

flowjo\_to\_gatingset( ws, name = NULL, subset = list(), execute = TRUE,
path = ``\,``, cytoset = NULL, backend\_dir = tempdir(), backend =
get\_default\_backend(), includeGates = TRUE, additional.keys =''\(TOT",
  additional.sampleID = FALSE,
  keywords = character(),
  keywords.source = "XML",
  keyword.ignore.case = FALSE,
  extend_val = 0,
  extend_to = -4000,
  channel.ignore.case = FALSE,
  leaf.bool = TRUE,
  include_empty_tree = FALSE,
  skip_faulty_gate = FALSE,
  compensation = NULL,
  transform = TRUE,
  fcs_file_extension = ".fcs",
  greedy_match = FALSE,
  mc.cores = 1,
  ...
)
In this case, I find that the “additional.keys” argument would likely work for this troubleshooting:
additional.keys
character vector: The keywords (parsed from FCS header) to be combined(concatenated with "_") with FCS filename to uniquely identify samples. Default is '\)TOT'
(total number of cells) and more keywords can be added to make this
GUID.

\begin{Shaded}
\begin{Highlighting}[]
\NormalTok{gs }\OtherTok{\textless{}{-}} \FunctionTok{flowjo\_to\_gatingset}\NormalTok{(}\AttributeTok{ws=}\NormalTok{ws, }\AttributeTok{name=}\DecValTok{1}\NormalTok{, }\AttributeTok{path =}\NormalTok{ Folder, }\AttributeTok{additional.keys=}\StringTok{"GROUPNAME"}\NormalTok{)}
\NormalTok{gs}
\end{Highlighting}
\end{Shaded}

A GatingSet with 6 samples

\begin{Shaded}
\begin{Highlighting}[]
\FunctionTok{class}\NormalTok{(gs)}
\end{Highlighting}
\end{Shaded}

{[}1{]} ``GatingSet'' attr(,``package'') {[}1{]} ``flowWorkspace''

Especially when working with CytoML, it is often good to have an idea of
how long it will take a particular function to run (to better plan how
to use our time while waiting, whether to go grab coffee, etc.). There
are a couple ways to do so.

One, using the system.time() function from base R, in which we surround
whatever line of code we wish to evaluate in \{\}

\begin{Shaded}
\begin{Highlighting}[]
\FunctionTok{system.time}\NormalTok{(\{}\FunctionTok{flowjo\_to\_gatingset}\NormalTok{(}\AttributeTok{ws=}\NormalTok{ws, }\AttributeTok{name=}\DecValTok{1}\NormalTok{, }\AttributeTok{path =}\NormalTok{ Folder, }\AttributeTok{additional.keys=}\StringTok{"GROUPNAME"}\NormalTok{)\})}
\end{Highlighting}
\end{Shaded}

user system elapsed 0.81 0.04 0.89

Alternatively, if we install the bench package, we can use the mark
function to evaluate how long it takes on average across numerous
iterations.

\begin{Shaded}
\begin{Highlighting}[]
\FunctionTok{install.packages}\NormalTok{(}\StringTok{"bench"}\NormalTok{)}
\end{Highlighting}
\end{Shaded}

\begin{Shaded}
\begin{Highlighting}[]
\FunctionTok{library}\NormalTok{(bench)}
\end{Highlighting}
\end{Shaded}

\begin{Shaded}
\begin{Highlighting}[]
\FunctionTok{mark}\NormalTok{(Test }\OtherTok{\textless{}{-}} \FunctionTok{flowjo\_to\_gatingset}\NormalTok{(}\AttributeTok{ws=}\NormalTok{ws, }\AttributeTok{name=}\DecValTok{1}\NormalTok{, }\AttributeTok{path =}\NormalTok{ Folder, }\AttributeTok{additional.keys=}\StringTok{"GROUPNAME"}\NormalTok{), }\AttributeTok{iterations=} \DecValTok{5}\NormalTok{)}
\end{Highlighting}
\end{Shaded}

\section{A tibble: 1 × 13}\label{a-tibble-1-13}

expression min median \texttt{itr/sec} mem\_alloc \texttt{gc/sec} n\_itr
n\_gc total\_time result memory time gc\\
\textless bch:\textgreater{} \\
1 ``Test \textless- flowjo\_to\_gatingset\ldots{} 859ms 865ms 1.15
13.3KB 0 5 0 4.35s

Now that we have loaded the contents of the FlowJo/Floreada workspace,
we can start exploring the various infrastructural capabilities of a
GatingSet object.

Let's start by evaluating whether the manually-drawn gates I drew
survived the journey. To do this, I can generate a visual gating treee
using the plot() function.

\begin{Shaded}
\begin{Highlighting}[]
\FunctionTok{plot}\NormalTok{(gs)}
\end{Highlighting}
\end{Shaded}

\pandocbounded{\includegraphics[keepaspectratio]{C:/Users/Macarena/MacarenaLM/notes/week04/04_IntroToTidyverse/week 05 Gating sets notes/imagenweek5.png}}




\end{document}
